%% ============================================================
%%  Chapitre 8 — Méta-learning cross-domaine
%% ============================================================
\chapter{Méta-learning cross-domaine}
\label{ch:metalearning}

Ce chapitre présente la contribution scientifique centrale du mémoire :
le méta-learning pour la généralisation cross-domaine du scorer de qualité.
Nous formalisons le problème, décrivons les meta-features extraites,
présentons les méta-learners évalués, et discutons les résultats de la
validation de l'hypothèse H3.

\section{Motivation et formalisation}
\label{sec:meta_motivation}

\subsection{Le problème de la généralisation}

Le Chapitre~\ref{ch:automl} a démontré que les pondérations optimales
du scorer varient significativement d'un domaine à l'autre (tableau~\ref{tab:weight_profiles}).
Cette observation soulève un problème pratique : lorsqu'un utilisateur
soumet un nouveau domaine de recherche, il est impossible de lancer une
optimisation Optuna complète (qui nécessite un oracle préalablement construit)
avant d'avoir produit les premiers résultats.

Deux approches naïves existent pour contourner ce problème :
\begin{enumerate}
  \item \textbf{Default weights} : utiliser des pondérations uniformes,
    ignorant toute connaissance accumulée sur les autres domaines.
  \item \textbf{Transfer Direct} : utiliser la moyenne des pondérations
    apprises sur tous les domaines connus, supposant qu'une pondération
    ``moyenne'' sera approximativement correcte pour tout nouveau domaine.
\end{enumerate}

Ces deux approches sont des baselines. Nous proposons une troisième voie :
\textbf{le méta-learning}, qui prédit les pondérations optimales d'un nouveau
domaine à partir de caractéristiques mesurables du corpus de ce domaine,
sans nécessiter ni oracle ni optimisation préalable.

\subsection{Formalisation}

Soient $\mathcal{D} = \{d_1, \ldots, d_N\}$ les $N = 19$ domaines d'évaluation.
Pour chaque domaine $d_i$, nous disposons de :
\begin{itemize}
  \item Un vecteur de meta-features $\phi(d_i) \in \mathbb{R}^M$ décrivant
    les caractéristiques statistiques du corpus de $d_i$ ;
  \item Les pondérations optimales $\mathbf{w}^*(d_i) \in \mathcal{W}$
    apprises par Optuna sur l'oracle de $d_i$.
\end{itemize}

Le méta-learner est une fonction $f_\theta : \mathbb{R}^M \to \mathcal{W}$
apprise sur $\{(\phi(d_i), \mathbf{w}^*(d_i))\}_{i=1}^N$.
Pour un nouveau domaine $d_{\text{new}}$, les pondérations prédites sont
$\hat{\mathbf{w}}(d_{\text{new}}) = f_\theta(\phi(d_{\text{new}}))$.

\section{Extraction des meta-features}
\label{sec:meta_features}

\subsection{Définition des meta-features}

Les meta-features sont des statistiques descriptives calculées directement
sur le corpus de publications d'un domaine, \textbf{sans utiliser l'oracle}
(pour éviter toute fuite d'information). Elles capturent des propriétés
comme la maturité du domaine, la distribution des citations, ou le niveau
d'expertise des auteurs.

Le tableau~\ref{tab:meta_features_list} liste l'ensemble des 13 features
calculées dans notre système.

\begin{table}[htbp]
  \centering
  \caption{Meta-features extraites pour chaque domaine d'évaluation.}
  \label{tab:meta_features_list}
  \rowcolors{2}{lightgray}{white}
  \begin{tabular}{llp{0.45\textwidth}}
    \toprule
    \textbf{Feature} & \textbf{Type} & \textbf{Description} \\
    \midrule
    \texttt{grade2\_ratio} & ratio & Ratio d'articles grade-2 dans l'oracle \\
    \texttt{mean\_h\_index} & continu & H-index moyen des auteurs principaux \\
    \texttt{citation\_median} & continu & Médiane des citations par article \\
    \texttt{citation\_gini} & continu & Coefficient de Gini des citations \\
    \texttt{pct\_high\_hindex} & ratio & \% d'auteurs avec h-index $>$ 20 \\
    \texttt{pct\_recent} & ratio & \% d'articles des 3 dernières années \\
    \texttt{pct\_high\_cited} & ratio & \% d'articles avec $>$ 100 citations \\
    \texttt{year\_std} & continu & Écart-type des années de publication \\
    \texttt{gold\_ratio} & ratio & Ratio grade-2 / corpus total \\
    \texttt{median\_year} & entier & Année médiane de publication \\
    \texttt{unique\_author\_ratio} & ratio & Ratio d'auteurs uniques / total auteurs \\
    \texttt{pct\_q1} & ratio & \% d'articles publiés dans des revues Q1 \\
    \texttt{topic\_concept\_overlap} & ratio & Recouvrement topic / concepts OpenAlex \\
    \bottomrule
  \end{tabular}
\end{table}

\subsection{Sélection des meta-features pertinentes}
\label{sec:feature_selection}

La sélection des features pertinentes est effectuée par analyse des
corrélations de Pearson entre chaque meta-feature et chaque dimension
de pondération, sur l'ensemble des $N = 19$ domaines. Le tableau~\ref{tab:pearson_correlations}
présente les corrélations maximales (en valeur absolue) obtenues.

\begin{table}[htbp]
  \centering
  \caption{Corrélations de Pearson (max $|r|$ sur les 6 pondérations), n=19 domaines.}
  \label{tab:pearson_correlations}
  \rowcolors{2}{lightgray}{white}
  \begin{tabular}{lccl}
    \toprule
    \textbf{Feature} & \textbf{max $|r|$} & \textbf{Sélectionnée ?}
      & \textbf{Corrélation principale} \\
    \midrule
    \texttt{grade2\_ratio}    & \textbf{0.67} & \checkmark & velocity ($r = -0.67$) \\
    \texttt{mean\_h\_index}   & \textbf{0.61} & \checkmark & relevance ($r = -0.61$) \\
    \texttt{citation\_median} & \textbf{0.58} & \checkmark & recency ($r = +0.58$) \\
    \texttt{citation\_gini}   & \textbf{0.56} & \checkmark & venue ($r = -0.56$) \\
    \texttt{pct\_high\_hindex}& \textbf{0.52} & \checkmark & relevance ($r = -0.52$) \\
    \texttt{pct\_recent}      & \textbf{0.45} & \checkmark & relevance ($r = -0.45$) \\
    \texttt{pct\_high\_cited} & \textbf{0.41} & \checkmark & recency ($r = +0.41$) \\
    \texttt{gold\_ratio}      & 0.39 & \texttimes & — \\
    \texttt{year\_std}        & 0.38 & \texttimes & — \\
    \texttt{unique\_author\_ratio} & 0.18 & \texttimes & — \\
    \bottomrule
  \end{tabular}
\end{table}

Les 7 features avec $\max|r| > 0.40$ constituent l'\texttt{EFFECTIVE\_FEATURES}
utilisé par le méta-learner, défini dans \texttt{src/metalearning/meta\_learner.py}.

\paragraph{Interprétation des corrélations.}
La corrélation la plus forte est celle entre \texttt{grade2\_ratio} et la
pondération velocity ($r = -0.67$) : les domaines avec une plus grande
proportion d'articles de référence (\textit{landmarks}) accordent moins
d'importance à la vélocité des citations, car les articles fondateurs restent
pertinents même s'ils ne sont plus activement cités récemment. La corrélation
entre \texttt{mean\_h\_index} et la pondération relevance ($r = -0.61$) indique
que les domaines à experts très reconnus (h-index élevé) ont moins besoin
de la pertinence sémantique : la réputation de l'auteur est suffisamment
discriminante.

\section{Méta-learners évalués}
\label{sec:meta_learners}

Nous évaluons 12 méta-learners, répartis en quatre familles.

\subsection{Méthodes à base de voisins (kNN)}

Le méta-learner kNN($k$) prédit les pondérations d'un domaine test $d$
comme une moyenne pondérée par la distance inverse des pondérations des
$k$ domaines d'entraînement les plus proches dans l'espace des meta-features :

\begin{equation}
  \hat{\mathbf{w}}_{\text{kNN}(k)}(d) \propto
  \sum_{i \in \mathcal{N}_k(d)} \frac{1}{\|\phi(d) - \phi(d_i)\|_2}
  \cdot \mathbf{w}^*(d_i)
  \label{eq:knn_prediction}
\end{equation}

où $\mathcal{N}_k(d)$ est l'ensemble des $k$ voisins les plus proches.
Les meta-features sont normalisées par MinMaxScaler avant le calcul de
distance, pour éviter que \texttt{citation\_median} (grande amplitude)
ne domine l'espace métrique.

Nous évaluons $k \in \{1, 2, 3\}$.

\subsection{Méthodes de régression linéaire}

Trois régresseurs linéaires sont évalués, chacun ajusté indépendamment
pour chaque dimension de pondération :

\paragraph{RidgeCV.} Régression ridge avec validation croisée automatique
pour la sélection du paramètre de régularisation $\alpha$ parmi
$\{0.01, 0.1, 1, 10, 100, 1000\}$.

\paragraph{BayesianRidge.} Régression bayésienne de crête qui ajuste
automatiquement les hyperparamètres de régularisation via l'inférence
bayésienne (maximisation de la log-vraisemblance marginale).

\paragraph{GPR (\textit{Gaussian Process Regression}).}
Régression par processus gaussien avec noyau RBF + constant + bruit blanc.
Cette méthode modélise l'incertitude de prédiction et est théoriquement
bien adaptée aux données de faible taille ($n = 18$ en LOO), mais souffre
en pratique d'instabilité numérique sur certains domaines.

\subsection{Méthodes d'ensemble (rank fusion)}

Ces méthodes combinent les prédictions des quatre base-learners
kNN(1), kNN(2), kNN(3) et BayesianRidge en un vecteur de pondérations
unique.

\paragraph{Soft Vote.} Moyenne arithmétique des prédictions des 4 membres,
suivie d'une normalisation :
\begin{equation}
  \hat{\mathbf{w}}_{\text{SoftVote}} = \frac{1}{4} \sum_{m=1}^{4} \hat{\mathbf{w}}_m
  \label{eq:soft_vote}
\end{equation}

\paragraph{Hard Vote (médiane).} Médiane par dimension sur les 4 membres
(avec 4 valeurs : moyenne des 2e et 3e valeurs ordonnées), suivie d'une
normalisation. Plus robuste aux prédictions aberrantes (\textit{outliers})
qu'une moyenne.

\paragraph{Borda Count.} Pour chaque membre $m$, les 6 dimensions de
pondération sont classées par valeur décroissante. La dimension en
position $r$ (sur 6) reçoit un score de Borda $b = 6 - 1 - r = 5 - r$.
Les scores sont sommés sur les 4 membres et normalisés :

\begin{equation}
  \hat{\mathbf{w}}_{\text{Borda}} \propto \sum_{m=1}^{4}
  \mathbf{b}(\hat{\mathbf{w}}_m)
  \label{eq:borda}
\end{equation}

où $\mathbf{b}(\mathbf{w})_j = 5 - \text{rank}_j(\mathbf{w})$.

\paragraph{Reciprocal Rank Fusion (RRF).} Score RRF par dimension :
\begin{equation}
  \text{RRF}_j = \sum_{m=1}^{4} \frac{1}{k + \text{rank}_j(\hat{\mathbf{w}}_m)}
  \label{eq:rrf}
\end{equation}
avec $k = 1$ (choix adapté à $n = 6$ dimensions : $k = 60$ standard
sur-lisserait les différences).

\subsection{Ensembles hybrides kNN--régression}

Deux ensembles convexes combinent kNN(1) avec un régresseur linéaire :

\begin{align}
  \hat{\mathbf{w}}_{\text{En+R}}(\alpha) &= \alpha \cdot \hat{\mathbf{w}}_{\text{kNN}(1)}
    + (1 - \alpha) \cdot \hat{\mathbf{w}}_{\text{Ridge}} \\
  \hat{\mathbf{w}}_{\text{En+B}}(\alpha) &= \alpha \cdot \hat{\mathbf{w}}_{\text{kNN}(1)}
    + (1 - \alpha) \cdot \hat{\mathbf{w}}_{\text{Bayes}}
\end{align}

Le paramètre $\alpha \in \{0.0, 0.1, \ldots, 1.0\}$ est sélectionné par
balayage sur le jeu d'entraînement LOO.

\section{Protocole d'évaluation : LOO-CV}

Chaque méta-learner est évalué par validation croisée Leave-One-Out sur
les 19 domaines. À chaque fold $i$ :
\begin{enumerate}
  \item Le méta-learner est ajusté sur les 18 domaines $\mathcal{D} \setminus \{d_i\}$ ;
  \item Il prédit $\hat{\mathbf{w}}(d_i)$ à partir de $\phi(d_i)$ ;
  \item Le NDCG@15 est calculé en classant les articles du corpus $C_{d_i}$
    avec les pondérations prédites, en utilisant les sous-scores précalculés
    pour accélérer le calcul.
\end{enumerate}

La performance d'un méta-learner sur un domaine est comparée au Transfer
Direct (pondérations moyennes sur les 18 domaines d'entraînement).
Un méta-learner remporte une ``victoire'' sur un domaine s'il surpasse
le Transfer Direct de plus de $10^{-6}$ en NDCG@15.

\section{Résultats : validation de H3}
\label{sec:h3_results}

\subsection{Taux de victoire par méthode}

Le tableau~\ref{tab:loo_results} présente les taux de victoire (vs Transfer
Direct) et les moyennes de NDCG@15 pour les 12 méthodes évaluées.

\begin{table}[htbp]
  \centering
  \caption{Résultats LOO-CV : victoires vs Transfer Direct et NDCG@15 moyen.
  Les méthodes dépassant le seuil de 75\,\% sont mises en évidence.}
  \label{tab:loo_results}
  \rowcolors{2}{lightgray}{white}
  \begin{tabular}{lccr}
    \toprule
    \textbf{Méthode} & \textbf{Wins/19} & \textbf{\%} & \textbf{Mean NDCG@15} \\
    \midrule
    \rowcolor{green!20}
    \textbf{Borda} & \textbf{15/19} & \textbf{79\,\%} $\checkmark$ & \textbf{0.282} \\
    RRF & 14/19 & 74\,\% & 0.254 \\
    En+R ($\alpha = 0.1$) & 13/19 & 68\,\% & 0.253 \\
    BayesRidge & 12/19 & 63\,\% & 0.253 \\
    En+B ($\alpha = 0.6$) & 12/19 & 63\,\% & 0.258 \\
    kNN(2) & 11/19 & 58\,\% & 0.249 \\
    kNN(3) & 11/19 & 58\,\% & 0.254 \\
    Ridge & 11/19 & 58\,\% & 0.249 \\
    SoftVote & 10/19 & 53\,\% & 0.252 \\
    HardVote & 9/19 & 47\,\% & 0.235 \\
    kNN(1) & 8/19 & 42\,\% & 0.259 \\
    GPR & 7/19 & 37\,\% & 0.235 \\
    \midrule
    \textit{Transfer Direct} & \textit{(référence)} & — & 0.247 \\
    \textit{Default weights} & \textit{(référence)} & — & 0.229 \\
    \textit{Optuna*} & \textit{(borne sup.)} & — & 0.387 \\
    \bottomrule
  \end{tabular}
  \small{*Optuna = pondérations réelles apprises sur le domaine test (non disponibles en pratique).}
\end{table}

\begin{resultbox}
  \textbf{H3 confirmée.} Borda Count atteint 15/19 = 79\,\% de taux de
  victoire sur le Transfer Direct, dépassant le seuil de 75\,\% fixé par H3.
  Sa moyenne de NDCG@15 est 0.282, soit $+14.2\,\%$ par rapport au Transfer
  Direct (0.247) et $+23.1\,\%$ par rapport aux pondérations par défaut (0.229).
\end{resultbox}

\subsection{Analyse par domaine}

La figure~\ref{fig:wins_heatmap} présente le tableau détaillé des gains
($\Delta$NDCG@15 vs Transfer Direct) par domaine et par méthode.
Les domaines sont triés par delta Borda décroissant.

\begin{figure}[htbp]
  \centering
  \includegraphics[width=\textwidth]{figures/diag_method_wins.png}
  \caption{Tableau des gains NDCG@15 (vs Transfer Direct) par domaine et méthode.
  Vert = victoire, rouge = défaite, ligne séparatrice après le 15\textsuperscript{e}
  domaine (médiane Borda).}
  \label{fig:wins_heatmap}
\end{figure}

Quatre domaines résistent à l'ensemble des méta-learners :
\textit{Transfer Learning} ($\Delta_{\text{Borda}} = -0.009$),
\textit{Fake News Detection} ($\Delta_{\text{Borda}} = -0.029$),
\textit{Speech ASR} ($\Delta_{\text{Borda}} = -0.107$), et
\textit{Anomaly Detection} ($\Delta_{\text{Borda}} = -0.148$).
Ces domaines présentent des profiles atypiques que les meta-features
actuelles ne capturent pas fidèlement, ce qui est discuté à la
section~\ref{sec:meta_discussion}.

\section{Test statistique de Wilcoxon}
\label{sec:wilcoxon}

\subsection{Configuration du test}

Afin d'évaluer si l'amélioration de Borda sur le Transfer Direct est
statistiquement significative et non due au hasard, nous appliquons
le test de Wilcoxon signé \cite{wilcoxon1945individual} (voir
Section~\ref{sec:etat_art_wilcoxon} pour la définition formelle) :

\begin{itemize}
  \item \textbf{Observations appariées} : $\Delta_i = \text{NDCG}_{Borda}(d_i) -
    \text{NDCG}_{Transfer}(d_i)$ pour $i = 1, \ldots, 19$ ;
  \item \textbf{Test unilatéral} : $H_1 : \text{median}(\Delta) > 0$ ;
  \item \textbf{Seuil de significativité} : $\alpha = 0.05$.
\end{itemize}

\subsection{Résultats}

\begin{table}[htbp]
  \centering
  \caption{Test de Wilcoxon signé (unilatéral) : Borda vs Transfer Direct, $n=19$.}
  \label{tab:wilcoxon_results}
  \begin{tabular}{lrrrll}
    \toprule
    \textbf{Méthode} & \textbf{$W^+$} & \textbf{$p$-value} & \textbf{$r_{rb}$}
      & \textbf{Effet} & \textbf{Sig. ($\alpha=0.05$) ?} \\
    \midrule
    \textbf{Borda} & \textbf{149.0} & \textbf{0.0145} & \textbf{+0.568}
      & \textbf{large} & \textbf{\checkmark OUI} \\
    En+R & 130.5 & 0.077 & +0.368 & medium & non \\
    RRF & 131.0 & 0.078 & +0.379 & medium & non \\
    BayesRidge & 127.5 & 0.095 & +0.345 & medium & non \\
    En+B & 126.0 & 0.113 & +0.326 & medium & non \\
    Ridge & 117.0 & 0.198 & +0.232 & small & non \\
    kNN(1) & 96.0 & 0.492 & +0.011 & small & non \\
    GPR & 87.5 & 0.619 & $-$0.143 & small & non \\
    \bottomrule
  \end{tabular}
\end{table}

\begin{resultbox}
  Borda Count est la seule méthode statistiquement significative à $\alpha = 0.05$
  ($W^+ = 149$, $p = 0.014$, $r_{rb} = +0.568$, effet large).
  En+R, RRF et BayesRidge montrent des effets medium ($r_{rb}$ entre 0.33 et 0.38)
  et des $p$-values entre 0.07 et 0.11, suggérant des tendances positives
  qui ne peuvent être confirmées avec $n = 19$.
\end{resultbox}

La figure~\ref{fig:wilcoxon_plot} illustre les $p$-values et tailles d'effet
pour les 12 méthodes.

\begin{figure}[htbp]
  \centering
  \includegraphics[width=\textwidth]{figures/phase8d_wilcoxon.png}
  \caption{$p$-values (haut, barres vertes = significatives) et corrélations
  rang-bisériales $r_{rb}$ (bas, effet size) pour les 12 méta-learners.
  La ligne orange tiretée indique $\alpha = 0.05$.}
  \label{fig:wilcoxon_plot}
\end{figure}

\section{Discussion}
\label{sec:meta_discussion}

\subsection{Pourquoi Borda excelle-t-il ?}

Le succès du Borda Count s'explique par sa robustesse aux prédictions
aberrantes. En agrégeant des rangs plutôt que des valeurs absolues,
il réduit l'impact des cas où un membre du comité (kNN ou BayesRidge)
produit une prédiction très éloignée de la réalité. Pour $n = 18$ domaines
d'entraînement, la variabilité des prédictions individuelles est élevée,
et les méthodes d'agrégation de rangs sont connues pour être plus robustes
en régime de faible données \cite{dwork2001rank}.

\subsection{Domaines difficiles}

Les quatre domaines où Borda échoue (\textit{Transfer Learning}, \textit{Fake
News}, \textit{Speech ASR}, \textit{Anomaly Detection}) partagent une
caractéristique commune : leurs profils de pondération optimaux sont
atypiques par rapport au reste du corpus. Par exemple, \textit{Anomaly
Detection} requiert des pondérations très concentrées sur l'impact
($w_{\text{impact}} = 0.39$) et l'auteur ($w_{\text{authors}} = 0.28$),
un profil non prédit par les meta-features actuelles.

Ces cas difficiles suggèrent que l'ensemble actuel de meta-features
n'est pas suffisamment expressif pour capturer toute la diversité des
domaines. Des meta-features supplémentaires — comme le ratio de preprints
arXiv, la densité sémantique du corpus, ou l'entropie des concepts
OpenAlex — pourraient améliorer la couverture des cas extrêmes.

\subsection{Limites et perspectives}

\paragraph{Taille du corpus d'évaluation.}
Avec $n = 19$ domaines, la puissance statistique du test de Wilcoxon est
limitée. Une étude avec $n \geq 30$ domaines permettrait de confirmer la
significativité d'En+R et RRF, actuellement à la frontière du seuil ($p \approx 0.08$).

\paragraph{Qualité des oracles.}
La construction des oracles par heuristique bibliographique (citations
dans des surveys) introduit un biais : les articles non cités dans les
surveys disponibles, même s'ils sont pertinents, sont classés grade-0.
Une annotation humaine sur un sous-ensemble permettrait de valider la
qualité de ces oracles.

\paragraph{Meta-features dynamiques.}
Les meta-features actuelles sont calculées une seule fois sur le corpus
initial. Des meta-features dynamiques, recalculées après chaque requête
de l'utilisateur, pourraient mieux capturer les propriétés spécifiques
de la requête plutôt que du domaine général.

\paragraph{Extension au méta-learning de niveau 2.}
La contribution actuelle correspond au Niveau 3 de généralisation
(méta-apprentissage). Une extension vers le Niveau 2 (fine-tuning rapide
sur quelques exemples du nouveau domaine) permettrait d'affiner les
prédictions lorsqu'un minimum de feedback utilisateur est disponible.
